\section{\ref{sec:glutcomputer}장. \nt{GLUT} 뉴런은 무엇을 계산하는가}

이 장의 논리적 명제는 음이 아닌 가중치로 된 회로에 관한 정리로, 장 안에서 유도했다. 실험은 뉴런 모델을 뒷받침하고, 이 정리들에 따라 조립한 회로(동시성 검출기, 지연선, 자기 유지 무리, 사슬)가 뇌에 실제로 있다는 것을 보여 준다.

{\footnotesize
\setlength{\tabcolsep}{3pt}\renewcommand{\arraystretch}{1.15}
\begin{longtable}{>{\raggedright}p{0.5cm}>{\raggedright}p{4.0cm}>{\raggedright}p{5.6cm}>{\raggedright}p{2.3cm}>{\raggedright\arraybackslash}p{1.7cm}}
\toprule
No. & 명제 & 실험과 측정한 것 & 출처 & 상태 \\
\midrule\endfirsthead
\toprule
No. & 명제 & 실험과 측정한 것 & 출처 & 상태 \\
\midrule\endhead
1 & 뉴런은 문턱값과 초기화가 있는 RC 회로이다~\eqref{eq:glutneuron} & 쥐 피질 피라미드 뉴런(절편)의 세포체에 살아 있는 뇌의 시냅스 흐름과 비슷한 잡음 전류를 주입했다. 평균 발화율이 전류의 평균과 산포에 어떻게 의존하는지는 느린 발화율 적응을 더한 누설 적분 발화 뉴런으로 기술된다. $\tau$와 지연의 범위는 장에서 실험 인용 없이 제시했다 & \cite{Rauch2003} & 측정됨 \\
2 & 모델~\eqref{eq:glutneuron}은 단순화이다. 입력 가지가 스스로 국소 스파이크를 낸다 & 생쥐 시각 피질 피라미드 뉴런의 가지돌기에서 in vivo로 직접 기록: 시각 자극이 NMDA 수용체에 의존하는 국소 가지돌기 스파이크를 일으키며, 이를 억누르면 뉴런의 방위 조율이 넓어진다 & \cite{Smith2013} & 측정됨 \\
3 & 일치 창 $\Delta_{\max}=\tau\ln\frac{w}{\theta-w}$~\eqref{eq:window}; $\theta=1.5w$이면 $0.7\tau$ & 모델~\eqref{eq:glutneuron}의 결과이다. 좁은 합산 창을 직접 측정한 자료는 빠른 억제가 있는 뉴런에 대해 있다(\ref{sec:gaba}장, 그 장 표의 3행) & 이 장의 유도 & 이론 \\
4 & \nt{GLUT} 뉴런만으로 된 네트워크는 입력의 단조 함수만 계산한다(명제~\ref{prop:monotone}) & 장에서의 증명: 우변이 감소하지 않는 침묵에서의 반복. 이것은 모델에 관한 명제이다. 억제 없는 살아 있는 네트워크에서 이를 검증한 실험은 없다 & 이 장의 유도 & 이론 \\
5 & 감각 기관은 전선 쌍을 준다. 어떤 세포는 자극이 강해질 때, 다른 세포는 약해질 때 응답한다 & 망막: 이미 쌍극 세포에서 원뿔 세포의 신호가 켜짐 채널과 꺼짐 채널로 갈라진다. 원숭이에서 켜짐 쌍극 세포를 약리학적으로 차단: 채널은 피질까지 따로 가며, 차단 뒤에는 빛이 강해지는 것은 잘 검출하지 못하지만 약해지는 것은 그렇지 않다 & \cite{Schiller1992} & 측정됨 \\
6 & 이중 레일 부호가 있으면 \nt{GLUT} 뉴런 네트워크는 공통 박자 없이 어떤 논리 함수든 비동기적으로 계산하며, 대가는 두 배의 뉴런 수이다(명제~\ref{prop:dualrail}, \eqref{eq:dualrail}, 그림~\ref{fig:dualxor}) & 이중 레일 부호와 신호 자체로 준비 상태를 판정하는 것은 지연에 둔감한 비동기 회로의 표준 기법이다. 뉴런 여섯 개짜리 배타적 OR 회로는 장에서 조립했다. 뇌가 논리 함수를 바로 이중 레일 부호로 계산한다는 것을 보여 주는 실험은 없다 & \cite{Sparso2001} & 이론 \\
7 & 사람에서 두 귀에 소리가 도착하는 시간의 최대 차는 $\approx0.6\ms$이고, 뇌는 약 10마이크로초를 구별한다 & 두 선택지 실험의 청취자: 시간 차 변별 문턱(정답 75\,\%)은 $150$--$1700$\,Hz 대역 잡음에서 $9$\,\textmu s, $1$\,kHz 순음에서 $11$\,\textmu s, 딸깍 소리에서 $28$\,\textmu s. 상한 $0.6\ms$는 장에서 $0.22\,\text{m}/343\,\text{m/s}$로 계산했다 & \cite{KlumppEady1956} & 측정됨 \\
8 & 새에서는 지연선과 동시성 검출기가 시간 차를 위치로 바꾼다~\eqref{eq:itd}(그림~\ref{fig:itd}) & 원숭이올빼미: 두 귀에서 온 축삭이 동시성 검출기 핵에 반대쪽에서 들어온다. 축삭을 따라 스파이크 도착 시간이 깊이에 따라 규칙적으로 변하고, 핵을 가로지르는 지연의 폭($700$\,\textmu m에 걸쳐 $160$\,\textmu s)은 두 귀 사이 지연의 범위와 일치한다 & \cite{Carr1990} & 측정됨 \\
9 & 포유류에서는 지연의 열이 발견되지 않았다. 같은 문제를 \nt{GLYC} 뉴런의 정확한 시간의 억제가 있는 동시성 검출기가 푼다 & 게르빌루스쥐 내측 상올리브에서 in vivo 기록: 응답이 방향 지도를 이루지 않는다. 미세 피펫으로 글리신과 그 차단제를 가하면, 시간을 정확히 맞춘 \emph{글리신} 억제가 지연의 작동 범위를 정한다는 것이 드러난다. 여기서 억제는 \nt{GABA}가 아니라 \nt{GLYC}에서 온다 & \cite{Brand2002,Grothe2010} & 측정됨 \\
10 & 되먹임이 강한 무리는 플립플롭이다. 자기 유지 활동은 자극 뒤 몇 초 동안 지속된다(그림~\ref{fig:bistable}) & 기억한 점으로 눈을 늦게 움직이는 과제(지연 $1$--$6$\,s)를 하는 원숭이의 전전두 피질: 과제와 관련된 뉴런 $170$개 중 $87$개가 지연 동안 발화율을 바꾸며, 그중 $79\,\%$에서 이 변화는 기억한 점의 방향에 의존한다. 이 활동을 두 안정 상태를 가진 되먹임이 유지한다는 것은 이론이다 & \cite{Funahashi1989,Wang2001} & 측정됨 \\
11 & 유입이 $\approx0.7\tau$보다 길면 비트가 써진다(그림~\ref{fig:recurrentgroup}b) & 방정식 $\tau\,\dd r/\dd t=-r+f(wr+I)$를 $w=1$, $I=0.6$에서 오일러 방법으로 계산했다. \texttt{models/}에 스크립트는 없다. 실험은 없다 & 이 장의 계산 & 모델 \\
12 & 기억은 스파이크 없이 거의 시냅스 촉진만으로 저장할 수 있다 & 이론: 말단의 잔류 칼슘이 기록을 약 1초 유지한다. 근거: 흰담비 내측 전전두 피질(여러 뉴런 동시 기록)에 여운이 긴 촉진성 시냅스로 연결된 피라미드 뉴런 하위 네트워크가 있다. 기억 유지 중 “조용한” 구간 관찰의 총설. 피질이 언제 어떤 방법을 쓰는지는 열린 문제이다 & \cite{Mongillo2008,WangMarkram2006,Stokes2015} & 가설 \\
13 & 기억은 피로로 지워진다. 신경전달물질 저장량이 고갈된다 & 피질 피라미드 뉴런의 쌍 기록: 스파이크가 잦으면 시냅스 응답이 줄고, 감소 속도는 방출 확률로 정해진다. 바로 이것이 자기 유지 무리를 끈다는 것은 직접 보이지 않았다 & \cite{Tsodyks1997} & 측정됨 \\
14 & 무리의 사슬은 사건이 시동하는 타이머이다. 명금류에서는 뉴런들이 엄격히 차례대로 발화한다 & 잠잘 때와 노래할 때 금화조의 고위 발성 중추에서 식별된 뉴런을 기록: 운동핵으로 출력을 보내는 뉴런은 저마다 노래의 정확한 한 순간에 짧은 버스트 하나를 내고, 뉴런들은 차례로 발화한다 & \cite{Hahnloser2002} & 측정됨 \\
\bottomrule
\end{longtable}}
